\subsection{Socioeconomic impacts}

\subsubsection{Economic impacts}

The economic impacts of \acp{bss} are multi-layered, extending from direct user savings to broader urban productivity and local economic development. While most \acp{bss} operate with public support or sponsorship and typically do not generate direct profits~\cite{Luo2019}, the true scale of their economic influence becomes apparent when externalities, such as time savings, health improvements, and local spending, are taken into account.

At the core of these impacts lie the day-to-day experiences of individual users. For many, bike-sharing reduces the monetary and temporal burdens of mobility, with users who transition from car, ride-hailing, or fare-based public transport benefitting from savings on costs like fuel, parking, vehicle upkeep, or transit fares. However, since a significant portion of bike-sharing trips replace walking or \ac{pt} rather than car use~\cite{Murphy2015, Fishman2013, Fishman2014, BachandMarleau2012}, the per-trip direct financial savings are often limited. Despite this, time savings can be meaningful, especially given the economic value assigned to leisure or saved commuting minutes~\cite{Qiu2018}.

When these individual time gains are aggregated, their cumulative effect becomes economically significant at the macro level. For example, research from Beijing shows that bike-sharing as a last-mile commuting solution reduces travel times by about 8 minutes per worker daily, translating to annual productivity gains between 0.3 and 1.2 billion yuan across the city, depending on scenario assumptions~\cite{Qiu2018}. Similarly, in Dublin, integration of bike-sharing through \textit{Dublinbikes} has led to an estimated annual saving of 16.5 million minutes for users compared to previous journeys—typically on foot or via \ac{pt}—amounting to approximately 6.06~million~\texteuro{} in direct time benefits~\cite{Bullock2017}. This increased accessibility also enhances effective urban density, driving an additional 6.79~million~\texteuro{} per year in wider agglomeration benefits~\cite{Bullock2017}. London provides comparable evidence, with the average cycle-hire trip being approximately 20\% shorter than the journey it replaces, further highlighting the efficiency advantages that \acp{bss} bring to dense cities~\cite{Woodcock2014}. Furthermore, these improvements in travel-time efficiency are not confined to Europe or China. In the U.S., for instance, empirical studies demonstrate that for many short \ac{od} journeys (under 3~km), bike-sharing is often as fast or even faster than taxis during rush hour congestion~\cite{FaghihImani2017_Hail}. Thus, across different urban contexts, bike-sharing consistently emerges as a tool for optimizing urban mobility.

Importantly, the economic contributions of \acp{bss} extend beyond time and cost savings to the stimulation of local commerce. By enhancing accessibility, bike-sharing enables trips that users might not otherwise make. This is supported by survey evidence from \textit{Nice Ride Minnesota}, where users reported making additional trips—primarily to retail and dining establishments—solely due to the convenience of bike-sharing~\cite{Schoner2012}. In just one week, respondents reported nearly \$35,000 in expenditures tied to these new or altered trips, averaging \$7–\$14 each. In Washington D.C., analysis of \textit{Capital Bikeshare} users showed that 16\% made trips they would not otherwise have taken and 23\% spent more money at destinations reached by bike than they would by other modes~\cite{Buehler2014}. Local businesses echoed these findings, with 20\% reporting increased sales attributable to nearby stations, 70\% perceiving positive neighborhood impacts, and 10\% noticing higher customer volumes~\cite{Buehler2014}. While some of this activity likely involves redistribution of spending rather than net new expenditure, the collective evidence points to bike-sharing as a driver of local economic vitality and commercial visibility.

The real estate sector represents yet another channel through which bike-sharing’s economic impacts are manifested. In Montreal, the 
introduction of the \ac{bss} was associated with higher property values for nearby multifamily 
homes. The analysis shows that each additional station within 800 meters raises the sale price 
of a unit by about \$709. Because the multifamily homes in the study area typically had around 
twelve stations within this buffer, this corresponds to an average increase of roughly \$8,650 
in sale price, or about 2.7\% of the mean value of units sold after the system was 
introduced~\cite{ElGeneidy2016}.

The real estate sector represents yet another channel through which bike-sharing’s economic impacts are manifested. In Montreal, the arrival of the \ac{bss} coincided with increased values for multifamily properties located near stations. Each new station within an 800-meter radius was associated with a \$709 rise in unit sale prices, and with a median proximity of twelve stations, this equated to roughly \$8,650 per unit—or approximately 2.7\% of the post-introduction average sale price~\cite{ElGeneidy2016}. This suggests benefits for property owners, with potential spillovers for city tax bases.

Health benefits add a further economic dimension. Increased physical activity through bike-sharing translates into reductions in premature mortality and morbidity, and these public health improvements are measurable in economic terms. For example, across 12 European cities, annual health benefits attributable to BSS use have been estimated at 18 million~\texteuro{}~\cite{Otero2018}; in the U.S., annual health-related economic gains have been valued at around \$36 million~\cite{Clockston2021}. More locally, a comprehensive cost–benefit analysis of \textit{Dublinbikes} found that benefits—notably health-related—outstripped costs, with benefit–cost ratios ranging from 1.21 to 1.68, ensuring a net societal gain for every euro invested~\cite{Murphy2015}.

Despite these diverse positive impacts, it is crucial not to overlook potential negative economic consequences. The uncontrolled boom of dockless \acp{bss} in China around 2017 exemplifies such risks. What initially appeared as rapid expansion quickly revealed itself as an unsustainable bubble, with some major operators experiencing mass layoffs and withdrawing from multiple cities~\cite{Fortune2017, BusinessInsider2018}. While these outcomes were driven largely by market forces, public agencies eventually bore the financial burden of dealing with abandoned bikes and urban clutter~\cite{Chen2019}. Thus, as cities leverage the many economic upsides of bike-sharing, prudent planning and regulation are needed to mitigate wider costs.



\subsubsection{Social equality}

While \acp{bss} are frequently promoted as instruments for advancing social equity~\cite{Shaheen2010}, their ability to do so is fundamentally tied to how effectively they reach populations with few mobility options. Proximity is critical: individuals living closer to stations are much more likely to use bike-sharing~\cite{Kabra2018, BachandMarleau2012}. As a result, inequitable spatial distribution of stations can unintentionally perpetuate or even deepen existing access inequalities. This dynamic is particularly pronounced for socially disadvantaged groups, who are often affected by structural transport poverty—marked by low car ownership, limited \ac{pt} access, and restricted connections to essential services such as employment, education, and healthcare~\cite{Lucas2016}. In these settings, if designed with equity in mind, bike-sharing can fill gaps by providing a flexible and affordable transport alternative~\cite{Shaheen2010}.

Extensive research highlights that spatial distribution is the main driver of inequalities in \acp{bss}~\cite{Teixeira2021, Goodman2014, Ogilvie2012}. When stations, service areas, and bicycles are concentrated in central, dense, and affluent neighborhoods, residents of peripheral or lower-income areas experience lower coverage and diminished access. Spatial fairness assessments across numerous cities show that bike-sharing provision systematically favors higher-density, higher-income districts~\cite{DuranRodas2020}. These station placement strategies are often aimed at maximizing ridership by targeting active social and economic hubs, however, they can inadvertently contribute to persistent socio-demographic imbalances in the user base (see Section~\ref{sec:who_uses_bss})~\cite{Ogilvie2012, Woodcock2014, Ricci2015}.

These spatial disparities are increasingly discussed within broader frameworks of transport equity and distributive justice~\cite{Cook2019, Litman2017, Pereira2017, Sheller2018}. Notably, even systems with free-floating bikes, which might theoretically overcome location-based barriers, continue to display inequities. For example, analyses in Shenzhen, China found that both the use and perceived accessibility of dockless bicycles are greater in wealthier, more central neighborhoods~\cite{Zhou2024}. Similarly, reviews of cycling equity more broadly reveal that infrastructure coverage, cycling activity, destination accessibility, and associated health benefits are all unevenly distributed, systematically favoring socially advantaged urban districts~\cite{Cunha2023}.

Beyond physical access, informational and awareness gaps constitute another layer of inequality. Even where equity-focused policies exist, disadvantaged residents may be less aware of \acp{bss} or targeted programs—often because of weaker social networks that include current users and more limited access to smartphones or the internet~\cite{McNeil2018}. This trend is evident in Montreal, where increased awareness of the \textit{BIXI} system over time was seen primarily among those living closer to stations and with higher education levels, while lower-education groups remained less informed~\cite{Bernatchez2015}.

Encouragingly, evidence shows that equitable outcomes are attainable when systems adopt inclusive policies and designs. For example, introducing \acp{e-b} alongside \acp{m-b} in Washington, D.C. broadened the demographic composition of users, drawing in more racially diverse, lower-income, and gender-balanced ridership~\cite{Bonilla-Alicea2020}. Additional studies demonstrate that interest in bike-sharing is often as high among disadvantaged groups as among advantaged ones~\cite{McNeil2018}, and sometimes per-capita usage even exceeds that of more privileged groups, provided that barriers to access are removed~\cite{Ogilvie2012}. This pattern suggests that disparities are less about user preferences and more about structural limitations. Therefore, expanding coverage into underserved areas, maintaining affordable pricing, and implementing targeted outreach and marketing can significantly enhance participation among marginalized communities~\cite{Goodman2014, McNeil2018, Ogilvie2012, Ursaki2015}.
